\documentclass[11pt]{article}

\usepackage[margin=1in]{geometry}
\usepackage[T1]{fontenc}
\usepackage[utf8]{inputenc}
\usepackage{lmodern}
\usepackage{microtype}
\usepackage{amsmath}
\usepackage{booktabs}
\usepackage{enumitem}
\usepackage{hyperref}
\hypersetup{colorlinks=true,linkcolor=blue,urlcolor=blue}

\title{Introducción al Proyecto --- Predicción de Churn en E-commerce (Olist)}
\author{
David Alan Delgado Zarco \\
Octavio Esquivel Alvarez del Castillo \\
John Williams Morrill \\
Jason Rosher Morrill
}
\date{\today}

\begin{document}
\maketitle
\tableofcontents
\newpage

% ------------------------------------------------------------------
\section{Motivación y pregunta de negocio}

El \textit{churn}, o abandono de clientes, representa uno de los mayores desafíos operativos en el comercio electrónico. Retener a un cliente existente es significativamente más
económico que adquirir uno nuevo; sin embargo, muchas plataformas carecen de mecanismos
automáticos para identificar, antes de que ocurra, qué clientes están en riesgo de no volver a
comprar.

Este proyecto aborda esa problemática sobre el \textbf{dataset público de Olist}, un marketplace
brasileño con datos reales de órdenes, clientes, pagos, reseñas e ítems. La pregunta central
que guía el trabajo es: \textbf{¿es posible anticipar qué clientes no volverán a comprar en los
próximos 90 días, a partir de su historial transaccional y de experiencia?}

Para responderla se construyó un pipeline completo de Machine Learning: desde la definición
rigurosa del label hasta el despliegue de un modelo en producción en AWS, incluyendo un
tablero interactivo de predicción individual.

% ------------------------------------------------------------------
\section{Datos y definición del label}

Se utilizaron cinco archivos del dataset de Olist: \texttt{orders}, \texttt{customers},
\texttt{payments}, \texttt{reviews} e \texttt{items}. Para evitar la duplicación de registros
provocada por relaciones 1:N, cada tabla auxiliar se agrega por \texttt{order\_id} antes de
unirse a la tabla de órdenes.

El aspecto más crítico del diseño fue construir el label de churn sin introducir
\textit{data leakage}. La solución adoptada es una \textbf{arquitectura de ventana temporal}:
se fija una \textbf{fecha de snapshot} y se observa si el cliente realiza alguna compra en la
\textbf{ventana futura} de $W = 90$ días siguiente. Formalmente:

\[
  \text{churn} = 1 \iff \text{el cliente no compra en } (\textit{snapshot},\ \textit{snapshot} + 90\text{ días}]
\]

Un cliente se etiqueta como $\text{churn}=0$ (no-churn / recompra) si sí compra dentro de esa
ventana. Adicionalmente, para excluir clientes históricamente inactivos que sesgarían el
modelo, la población se limita a aquellos con al menos una compra dentro de un
\textit{lookback} de 180 días previos al snapshot.

% ------------------------------------------------------------------
\section{Ingeniería de features y partición temporal}

Las variables se construyen a nivel de \texttt{customer\_unique\_id}, consolidando todas las
órdenes anteriores al snapshot en tres grupos:

\begin{itemize}[leftmargin=1.5em]
  \item \textbf{RFM:} \textit{Recency} (días desde la última compra), \textit{Frequency}
        (número de órdenes únicas) y \textit{Monetary} (suma total de pagos).
  \item \textbf{Experiencia y logística:} promedio del puntaje de review, días de entrega,
        días de retraso frente a la fecha estimada, número de ítems por orden, suma promedio
        de precio y de flete.
  \item \textbf{Contexto geográfico:} estado del cliente en Brasil
        (\texttt{customer\_state}), tratado como variable categórica con codificación
        \textit{one-hot}.
\end{itemize}

La partición train/test se realiza de forma \textbf{temporal} —no aleatoria— ordenando los
clientes por su última compra previa al snapshot y asignando el 80\% más antiguo a
entrenamiento y el 20\% más reciente a prueba. Esto reproduce un escenario realista de
generalización y elimina cualquier filtración de información futura.

% ------------------------------------------------------------------
\section{Modelos y evaluación}

Se entrenaron y compararon dos familias de modelos, ambas con corrección por desbalance de
clases:

\begin{itemize}[leftmargin=1.5em]
  \item \textbf{Regresión Logística} con \texttt{class\_weight=balanced} y regularización $L_2$.
  \item \textbf{Random Forest} con \texttt{class\_weight=balanced\_subsample} y
        \texttt{n\_estimators=400}.
\end{itemize}

Dado el fuerte desbalance del dataset —la mayoría de los clientes de Olist son compradores
únicos que no regresan—, la métrica principal adoptada es el
\textbf{PR-AUC para la clase no-churn} (\texttt{pr\_auc\_nochurn\_0}), que evalúa la
capacidad del modelo de identificar correctamente la recompra. El accuracy convencional
resulta engañoso en este contexto: un clasificador trivial que prediga siempre churn alcanza
alta exactitud pero fracasa completamente en detectar la clase de interés.

Un segundo aspecto clave fue el \textbf{ajuste del umbral de decisión}. Mediante un barrido
sobre 99 valores entre 0 y 1, se identificó el umbral que maximiza el $F_1$ para no-churn.
El valor resultante fue $\hat\tau = 0.12$, considerablemente inferior al 0.5 convencional,
lo cual refleja la rareza del evento de recompra en la población.

El modelo seleccionado fue la \textbf{Regresión Logística con $C=0.1$} y umbral operativo
$0.12$, registrado en MLflow como \texttt{logreg\_C0.1\_thr012} por presentar el mayor
\texttt{pr\_auc\_nochurn\_0} entre todos los runs comparados.

% ------------------------------------------------------------------
\section{Gestión de experimentos con MLflow en AWS EC2}

Todo el ciclo experimental se gestionó con \textbf{MLflow Tracking Server}, desplegado en
una instancia \textbf{EC2 Ubuntu} de AWS. El servidor almacena los parámetros, métricas y
artefactos —modelo serializado, matrices de confusión, curvas Precision-Recall— de cada
experimento, permitiendo su comparación visual en la interfaz web accesible por IP pública.

Esta infraestructura garantiza la reproducibilidad de los experimentos y la trazabilidad
completa del proceso de selección del modelo, tal como se esperaría en un entorno de
producción real.

% ------------------------------------------------------------------
\section{Tablero de predicción individual}

El proyecto culmina con el despliegue de un \textbf{tablero interactivo} accesible desde el
navegador. Su arquitectura se compone de tres capas:

\begin{itemize}[leftmargin=1.5em]
  \item \textbf{Frontend React/Vite:} captura las variables del cliente a través de un
        formulario lateral y muestra el resultado de la predicción.
  \item \textbf{API Express (Node.js):} recibe la solicitud en el endpoint
        \texttt{POST /api/predict} y orquesta la ejecución del proceso de inferencia.
  \item \textbf{Script de inferencia Python:} carga el modelo \texttt{best\_model\_logreg.joblib}
        y ejecuta \texttt{predict\_proba}, devolviendo la probabilidad de churn como JSON.
\end{itemize}

El flujo de datos de extremo a extremo es:

\[
  \text{Usuario} \xrightarrow{\text{POST}} \text{React/Vite} \rightarrow \text{Express}
  \xrightarrow{\texttt{spawn}} \text{Python} \rightarrow \texttt{joblib} \rightarrow
  \texttt{churnProbability}
\]

El tablero muestra la probabilidad de churn en porcentaje junto con el nivel de riesgo
(bajo, medio o alto) y, de forma destacada, un panel de \textbf{contribución de features}
que desglosa qué variables impulsaron la predicción hacia churn o hacia retención. Esta
interpretabilidad —derivada directamente de los coeficientes de la regresión logística—
convierte el resultado en una herramienta accionable para equipos de retención.

% ------------------------------------------------------------------
\section{Conclusiones}

El proyecto demuestra que es posible construir un sistema de predicción de churn funcional
y desplegado en producción con un stack accesible: datos abiertos de e-commerce, pipeline
reproducible con MLflow, modelo lineal interpretable y tablero en la nube.

Los aprendizajes técnicos centrales son:
\begin{itemize}[leftmargin=1.5em, itemsep=2pt]
  \item El diseño correcto del label mediante ventana temporal es indispensable para evitar
        \textit{data leakage} y obtener estimaciones de generalización válidas.
  \item En datasets desbalanceados, el ajuste del umbral de decisión tiene un impacto mayor
        que la elección del algoritmo.
  \item MLflow proporciona trazabilidad y reproducibilidad sin sobrecarga operativa
        significativa.
\end{itemize}

Como trabajo futuro se identifican: incorporar señales de comportamiento por categoría de
producto, explorar modelos de gradient boosting calibrados, y añadir monitoreo de deriva
del modelo (\textit{data drift}) en producción.

\end{document}
